% Krishnan's default LaTeX document style. Copy this file, then replace the title and body.
\documentclass[11pt,a4paper]{article}
\usepackage[margin=2cm]{geometry}
\usepackage{amsmath,amssymb}
\usepackage{tikz}
\usepackage{pgfplots}
\pgfplotsset{compat=1.18}
\usepackage{booktabs}
\usepackage{enumitem}
\usepackage{parskip}
\usepackage[pdfauthor={Krishnan},pdftitle={SEE + SEE = AXES}]{hyperref}

\newcommand{\colsum}[3]{%
\begin{center}\large
\renewcommand{\arraystretch}{1.3}
\begin{tabular}{c@{\hspace{1.2em}}c@{\hspace{1.2em}}c@{\hspace{1.2em}}c@{\hspace{1.2em}}c}
#1\\
#2\\ \midrule
#3
\end{tabular}
\end{center}}

\begin{document}

\begin{center}{\Large\bfseries SEE + SEE = AXES}\end{center}

\section*{Problem}

In the addition sum below, each letter stands for a different non-zero digit. What digit does $X$ stand for?

\colsum{ & & $S$ & $E$ & $E$}{$+$ & & $S$ & $E$ & $E$}{ & $A$ & $X$ & $E$ & $S$}

\section*{Solution}

Work through the columns from right to left, as in ordinary addition. When a column adds up to 10 or more, write its last digit (the \emph{units digit}) and pass 1 to the next column on the left. That 1 is the \emph{carry}. Adding two numbers, the carry is always 0 or 1.

\textbf{\boldmath Step 1: the digit $A$.} Both rows are the same number, so the sum is $2 \times SEE$. Since $SEE \le 999$, we get $2 \times SEE \le 1998$, which is less than 2000.

\textbf{\boldmath The 4-digit answer starts with 1, so $A = 1$.}

\textbf{\boldmath Step 2: the digit $E$ (tens column).} In the tens column, $E + E + \text{carry}$ has units digit $E$. Take one $E$ away: $E + \text{carry}$ has units digit 0. But $E + \text{carry}$ lies between 1 and 10, so it must equal 10. With carry 0 we would need $E = 10$, which is not a digit. So the carry is 1 and $E = 9$.

\textbf{\boldmath $E = 9$, and the units column must send a carry of 1.}

\textbf{\boldmath Step 3: the digit $S$ (units column).} $E + E = 9 + 9 = 18$. Write 8, carry 1 --- exactly the carry Step 2 needed.

\textbf{\boldmath $S = 8$.}

\textbf{\boldmath Step 4: the digit $X$ (hundreds column).} The tens column is $9 + 9 + 1 = 19$: write 9, carry 1. The hundreds column is $S + S + 1 = 8 + 8 + 1 = 17$: write 7, carry 1. That last carry is the leading 1, which is $A$.

\textbf{\boldmath $X = 7$.}

\section*{Check}

\colsum{ & & 8 & 9 & 9}{$+$ & & 8 & 9 & 9}{ & 1 & 7 & 9 & 8}

$899 + 899 = 1798$. The digits $S = 8$, $E = 9$, $A = 1$, $X = 7$ are all different and non-zero.

\begin{center}\textbf{\boldmath Answer: $X = 7$}\end{center}

\end{document}
